\documentclass[border=14pt,varwidth=18.5cm]{standalone}
\usepackage{amsmath,amssymb}
\usepackage{mathtools}
\usepackage{empheq}
\usepackage[T1]{fontenc}
\usepackage{lmodern}
\usepackage{xcolor}
\usepackage{booktabs}
\usepackage{array}
\usepackage[american,siunitx]{circuitikz}
\definecolor{head}{HTML}{0072B2}
\definecolor{hl}{HTML}{FFF59D}
\definecolor{note}{HTML}{555555}
\newcommand{\sheettitle}[1]{{\Large\bfseries\color{head}#1}\par\smallskip
  {\small\color{note}Jeremia Christ Immanuel Manalu - 5023231017 - Biocybernetics (EB234701), Class Q}\par\medskip}
\newcommand{\sect}[1]{\par\medskip{\large\bfseries\color{head}#1}\par\smallskip}
\newcommand{\hlbox}[1]{\colorbox{hl}{$\displaystyle #1$}}
\newcommand{\nomtable}[1]{\par\medskip{\large\bfseries\color{head}Notation (in order of appearance)}\par\smallskip
  \renewcommand{\arraystretch}{1.15}%
  \begin{tabular}{>{$}l<{$} p{10.6cm} l}
  \toprule
  \textrm{\textbf{Symbol}} & \textbf{Meaning} & \textbf{Unit} \\
  \midrule
  #1
  \bottomrule
  \end{tabular}}
\setlength{\parindent}{0pt}

\renewcommand{\theequation}{4.\arabic{equation}}
\begin{document}
\sheettitle{Cable Equation Derivation for Action Potential Propagation}

\sect{Discretized axon (RC transmission line, lecturer week 9)}
\begin{center}
\begin{circuitikz}[scale=0.82, transform shape]
  % outside rail
  \draw (-3.3,3.4) -- (0.8,3.4) to[R, l=$r_o\Delta x$] (4.2,3.4) -- (5.8,3.4) to[R, l=$r_o\Delta x$] (9.2,3.4) -- (10.8,3.4) to[R, l=$r_o\Delta x$] (14.2,3.4) -- (15.2,3.4);
  % inside rail
  \draw (-3.3,0) -- (0.8,0) to[R, l_=$r_i\Delta x$] (4.2,0) -- (5.8,0) to[R, l_=$r_i\Delta x$] (9.2,0) -- (10.8,0) to[R, l_=$r_i\Delta x$] (14.2,0) -- (15.2,0);
  % membrane elements at x, x+dx, x+2dx
  \foreach \xx/\lab in {0/x, 5/x{+}\Delta x, 10/x{+}2\Delta x}{
    \draw (\xx-0.8,3.4) to[C, l_=$c_m\Delta x$] (\xx-0.8,0);
    \draw (\xx+0.8,3.4) to[R, l=$g_m\Delta x$] (\xx+0.8,0);
    \node at (\xx,-0.8) {$\lab$};
  }
  \draw[->, thick] (1.7,0.6) -- node[above]{$i(x)$} (3.3,0.6);
  \draw[->, thick] (6.7,0.6) -- node[above]{$i(x{+}\Delta x)$} (8.3,0.6);
  \draw[->, thick] (-2.9,2.6) -- node[left]{$i_m(x)$} (-2.9,0.8);
  \node at (6.3,4.45) {\textbf{Outside}};
  \node at (6.3,-1.45) {\textbf{Inside}};
\end{circuitikz}
\end{center}
In the active (HH) cable every membrane element $(c_m,\ g_m)$ is replaced by the HH circuit of sheet 1.

\sect{Step 1 - Ohm's law along the axis}
\begin{align}
v(x+\Delta x,t)-v(x,t) &= -\,i(x,t)\,(r_o+r_i)\,\Delta x \\
\lim_{\Delta x\to 0}\ \frac{v(x+\Delta x)-v(x)}{\Delta x} &= \frac{\partial v}{\partial x} = -\,(r_o+r_i)\,i
\end{align}

\sect{Step 2 - Kirchhoff's current law at a node}
\begin{align}
i(x+\Delta x,t)-i(x,t) &= -\,i_m(x,t)\,\Delta x,\qquad i_m = c_m\frac{\partial v}{\partial t}+i_{ion} \\
\frac{\partial i}{\partial x} &= -\Big(c_m\frac{\partial v}{\partial t}+i_{ion}\Big)
\end{align}

\sect{Step 3 - eliminate the axial current}
Differentiating (4.2) with respect to $x$ and substituting (4.4):
\begin{equation}
\hlbox{\frac{\partial^2 v}{\partial x^2} = (r_o+r_i)\Big(c_m\frac{\partial v}{\partial t}+i_{ion}\Big)}
\end{equation}

\sect{Step 4 - explicit form used by the program}
\begin{equation}
\hlbox{\frac{\partial v}{\partial t} = D\,\frac{\partial^2 v}{\partial x^2} - \frac{i_{ion}}{c_m}
= D\,\frac{\partial^2 v}{\partial x^2} + \frac{J_{in}-J_{ion}}{C_m},\qquad D=\frac{1}{(r_o+r_i)\,c_m}}
\end{equation}
\begin{align}
r_i &= \frac{R_i}{\pi a^2},\qquad c_m = 2\pi a\,C_m,\qquad i_{ion} = 2\pi a\,J_{ion} \\
J_{ion} &= \begin{cases} J_{Na}+J_K+J_L & \text{active HH cable (propagating AP)}\\ g_m\,v & \text{passive RC cable (electrotonic decay)}\end{cases}
\end{align}
The lecture slide writes the coefficient as $\big[(r_o+r_i)\,C_m\big]^{-1}$ with $C_m$ the
capacitance per unit length (the $c_m$ of this sheet, e.g.\ $1.6\times10^{-9}$ mF/mm), which is the same $D$. For the passive cable the steady spread decays as $e^{-x/\lambda}$ with
$\lambda=\sqrt{1/\big(2\pi a\,g_m\,(r_o+r_i)\big)}$.

\nomtable{
x,\ \Delta x & position along the axon and segment length (displayed in mm by the program) & cm \\
r_o,\ r_i & external and internal (axial) resistance per unit length & $\Omega$/cm \\
i & axial current & $\mu$A \\
c_m & membrane capacitance per unit length, $2\pi a C_m$ & $\mu$F/cm \\
g_m & membrane conductance of the passive cable & mS/cm$^2$ \\
i_m & membrane current per unit length & $\mu$A/cm \\
v & membrane potential displacement from rest & mV \\
t & time & ms \\
i_{ion} & ionic membrane current per unit length & $\mu$A/cm \\
D & diffusion coefficient of the cable & cm$^2$/ms \\
J_{in},\ J_{ion} & stimulus and ionic current densities & $\mu$A/cm$^2$ \\
C_m & membrane capacitance per unit area & $\mu$F/cm$^2$ \\
R_i & specific resistivity of the axoplasm & $\Omega\,$cm \\
a & axon radius & cm \\
J_{Na},J_K,J_L & HH ionic current densities (sheet 1) & $\mu$A/cm$^2$ \\
\lambda & passive length constant & cm \\
}
\end{document}
